% 封面 + 目录 + 导读
\begin{titlepage}
\vspace*{5cm}
\begin{center}
{\Huge\bfseries 算法竞赛风格题入门}\par\vspace{10pt}
{\Large LeetCode 上的竞赛味题目：从 L3 套路往上走一层}\par\vspace{24pt}
{\large 快速幂与取模 · 组合计数 · 图的形状 · 判环 · 函数图 · 建图 · 状压 DP · 轮廓线 DP · 搜索剪枝 · 构造}\par
\vfill
{\color{note} 2026-11-20 整理自 blog/docs/advanced 与 tracker 笔记\\
每章同一个组织方式：总览 $\to$ 构件（直觉 / 规则 / 例子）$\to$ 模板 $\to$ Usage 表 $\to$ 坑}
\end{center}
\end{titlepage}

\tableofcontents
\clearpage

\section*{导读}
\addcontentsline{toc}{section}{导读}

L3 的面试套路熟了以后，再往上走一层。这里的题都来自看板上的 \textbf{ACM 进阶}题单，是 LeetCode 上竞赛味比较重的题。
每章还是 L3 那个形状（怎么认 $\to$ 模板 $\to$ 栽过的坑），只是认题的信号变了。

\subsection*{和面试题的核心区别}

\begin{tabularx}{\linewidth}{@{}>{\bfseries}l X X@{}}
\toprule
& \textbf{L3 面试题} & \textbf{竞赛题}\\
\midrule
第一眼看什么 & 题面：给了树 / 区间 / 字符串 & \textbf{数据范围}：$n$ 多大，直接决定能用什么复杂度\\
题面和结构 & 题面直接告诉你是什么结构 & 题面是个故事，\textbf{先翻译成}图、函数图或计数模型，才轮得到套路\\
答案 & 构造出答案本身（一条路径、一个子数组） & 常常「对 $10^9+7$ 取模」：答案大到没法枚举，只能\textbf{数}出来\\
结构 & NeetCode 那十几类 & 多出几种面试基本见不到的：函数图、轮廓线、组合恒等式\\
构造题 & 几乎没有 & 有：没有搜索空间，靠观察直接给出一个合法答案\\
对错 & 思路对、能讲清楚就有分 & 只有 AC 和不 AC，边界、溢出、常数（Python 尤其）都得对\\
\bottomrule
\end{tabularx}

\subsection*{1. 数据范围就是提示}

竞赛题的数据范围是\textbf{故意卡的}，基本可以从它倒推出要用的算法：

\begin{center}
\begin{tabular}{@{}l l l l@{}}
\toprule
\textbf{$n$ 的上限} & \textbf{能跑的复杂度} & \textbf{通常意味着} & \textbf{本书哪章}\\
\midrule
$\le 10 \sim 12$ & $O(n!)$ & 全排列暴搜 & 第 8 章\\
$\le 15 \sim 20$ & $O(2^n \cdot n)$ & \textbf{状压}：集合当成一个整数 & 第 6 章\\
某一维 $\le 5$ & $O(\text{长维} \cdot 3^m \cdots)$ & 把小的那维压成状态 & 第 7 章\\
$\le 500$ & $O(n^3)$ & 区间 DP、Floyd & —\\
$\le 5000$ & $O(n^2)$ & 二维 DP、杨辉三角 & 第 1 章\\
$\le 2 \cdot 10^5$ & $O(n \log n)$ & 排序、堆、二分、线段树 & —\\
$\le 10^9$ & $O(\sqrt n)$ 或 $O(\log n)$ & 二分答案、数学 & —\\
$\le 10^{18}$ & $O(\log n)$ & \textbf{快速幂}、矩阵快速幂、倍增 & 第 1、2 章\\
\bottomrule
\end{tabular}
\end{center}

\begin{principle}
面试题也给数据范围，但很少靠它出题。到了这一层，\textbf{先看数据范围再读题}。
\end{principle}

\subsection*{2. 「取模」等于「这是计数题」}

题目要求答案取模，说明答案是指数级的，不可能一个一个列出来。能走的路只有两条：
组合数学直接算，或者计数 DP 一层一层数。模意义下加、减、乘都可以随时取模，
\textbf{除法不行}，要换成乘逆元。这一套在第 1 章。

\subsection*{3. 建模比套模板难}

L3 的难点在于「认出是哪个套路」，这一层的难点在\textbf{前面那一步}：
把「石头」「员工」「颜色」翻译成点、边和状态。
翻译完通常就落回 L3 的东西了（并查集、拓扑排序、DP），所以每章都会标它建在哪个 L3 套路上。

\subsection*{全书路线图}

\begin{tabularx}{\linewidth}{@{}c >{\bfseries}l X l l@{}}
\toprule
\textbf{章} & \textbf{主题} & \textbf{认题信号} & \textbf{建在 L3 的} & \textbf{代表题}\\
\midrule
1 & 快速幂、取模、组合计数 & 「对 $10^9+7$ 取模」「有多少种方式」；指数非常大 & 位运算、计数 DP & 372、2400\\
2 & 图的形状（图论总览） & 题目和环有关：先按方向 / 出度选支，再按每块的边数或有向环定形状 & DFS、BFS、并查集 & 思维导图\\
3 & 判环 & 「是否有环」「最短环多长」「哪些点永远走不进环」 & DFS、BFS & 802\\
4 & 函数图 & 输入是 \code{edges[i]} / \code{favorite[i]}：每个点恰好一条出边 & DFS、图 & 2127\\
5 & 建图：行和列当点 & 题面是网格、石头、约束，问的却是连通 / 先后顺序 & 并查集、拓扑 & 947\\
6 & 状压 DP & $n \le 20$；状态是「一个子集」 & 位运算、一维 DP & 1986\\
7 & 铺砖 / 轮廓线 DP & 网格铺砖、涂色，有一维 $\le 5$ & 状压 DP、二维 DP & 1931\\
8 & 搜索 + 剪枝 & 数据小得反常，想不出 DP 状态，贪心有反例 & 回溯 & 1240\\
9 & 构造 & 「返回任意一个满足条件的数组」 & 靠观察 & 932\\
\bottomrule
\end{tabularx}

\subsection*{怎么读}
\begin{itemize}
\item 每章开头的\textbf{总览图}先看：它说明这一章的工具谁是谁的前置，以及题目分成哪几种「化法」。
\item 构件按「直觉 $\to$ 规则 $\to$ 例子」写；蓝框是可以迁移到别的题的原则。
\item 每章的 \textbf{Usage 表}是题单：题 $\to$ 化成什么 $\to$ 用到哪个构件。时间不够就只做代表题。
\item 灰框「你踩过的坑」来自 tracker 里自己做题的笔记，只有做过的题才有。
\end{itemize}
